% DOC TYPE E DEF QB SOFTWARE %%%%%%%%%%%%%%%%%%%%%%%%%
\documentclass[12pt]{article}
%%%%%%%%%%%%%%%%%%%% PACKAGES %%%%%%%%%%%%%%%%%%%%
\usepackage[english,italian]{babel}
\usepackage[a4paper, margin=3cm]{geometry}
\usepackage[T1]{fontenc}
\usepackage[utf8]{inputenc}
\usepackage[table]{xcolor}
\usepackage{amsmath}
\usepackage{graphicx}
\usepackage{float}
\usepackage{tabularx}
\usepackage{booktabs}
\usepackage{hyperref}
\usepackage{xcolor}
\usepackage{multicol}
\usepackage{multirow}
\usepackage{soul}
\usepackage{enumitem}
\usepackage{textcomp}
\usepackage{eurosym}
\usepackage{lastpage}
\usepackage{fancyhdr}
\usepackage{adjustbox}
\usepackage{subfiles}

%%%%%%%%%%%%%%%%%%%% COMMAND -> New commands %%%%%%%%%%%%%%%%%%%%
\newcommand{\mailtoQBS}
{
	\href{mailto:qbsoftware.swe@gmail.com}{qbsoftware.swe@gmail.com}
}

\let\oldpar\paragraph
\renewcommand{\paragraph}[1]{\oldpar{#1}\mbox{}\\}

%%%%%%%%%%%%%%%%%%%% COMMAND -> Redefine LaTeX Commands %%%%%%%%%%%%%%%%%%%%
\def\title#1{\gdef\THETITLE{#1}}
\def\date#1{\gdef\THEDATE{#1}}
\def\footname#1{\gdef\THEFOOTNAME{#1}}

\let\oldtexteuro\texteuro
\renewcommand{\texteuro}{\euro}

%%%%%%%%%%%%%%%%%%%% STYLES -> Commands %%%%%%%%%%%%%%%%%%%%
\newcommand{\makefirstpage}
{
	\begin{titlepage}
		
		% Defines a new command for the horizontal lines, change thickness here
		\newcommand{\HRule}{\rule{\linewidth}{0.2mm}} 
		
		\center % Center everything on the page
		
		
		%	Heading Sections
		
		\textsc{\LARGE QB Software}\\[.1cm] 
		\includegraphics[scale=.15]{imgs/qb-software-logo.png}\\[-.1 cm] 
		\includegraphics[scale=.025]{imgs/x.png}\\[0.5cm]
		\includegraphics[scale=.3]{imgs/unipd_logo.png}\\[.5cm]
		\textsc{\Large Università degli studi di Padova}\\[0.5cm] 
		\textsc{\large corso di ingegneria del software }\\[0.5cm]
		\textsc{\large anno accademico 2023/2024 }\\[0.5cm]
		
		
		%	Title section and date section
		
		\ifdefined\THEDATE
		\HRule \\[0.4cm]
		{ \huge{ \bfseries {\THETITLE}} \\ [.5cm]
			\THEDATE}\\[0.4cm] 
		\HRule \\[0.4cm]
		\else
		\HRule \\[0.4cm]
		{ \huge{ \bfseries {\THETITLE}} } \\
		\HRule \\[0.4cm]
		\fi
		%
		\textsc{Contatti:} \mailtoQBS\\[0.3cm]
		
		\vfill 
		
	\end{titlepage}
}

%%%%%%%%%%%%%%%%%%%% ACCESSIBILITY %%%%%%%%%%%%%%%%%%%%
\let\oldhref\href
\renewcommand{\href}[2]{\oldhref{#1}{\textcolor{blue}{\ul{#2}}}}

\hypersetup{
	colorlinks = true,
	linkcolor = cyan,
}

%%%%%%%%%%%%%%%%%%%% STYLES -> Settings %%%%%%%%%%%%%%%%%%%%
% Enable header and footer style
\pagestyle{fancy}
\thispagestyle{empty}
\thispagestyle{fancy}
\pagestyle{fancy}

% Header
\setlength{\topmargin}{-40pt}
\setlength{\headsep}{60pt}
\fancyhf{}
\lhead{QB Software}
\setlength{\headheight}{15pt}
\rhead{\includegraphics[width=1cm]{imgs/qb-software-logo.png}}

% Footer 
\fancyfoot{}
\fancyfoot[L]{\THEFOOTNAME}
\fancyfoot[R]{Pagina \thepage~di~\pageref{LastPage}}
\futurelet\TMPfootrule\def\footrule{\TMPfootrule}
\setcounter{page}{0}
\pagenumbering{arabic}
\renewcommand{\footrulewidth}{0.3pt}

%%%%%%%%%%%%%%%%%%%% ENVIRONMENT %%%%%%%%%%%%%%%%%%%%
% Remove colorless padding in booktabs
\setlength{\aboverulesep}{0cm}
\setlength{\belowrulesep}{0cm}
\setlength{\extrarowheight}{.75ex}

\newenvironment{todo}{
	\rowcolors{2}{cyan!80!black!30!}{cyan!80!black!20!}
	\begin{tabular}{p{3.48cm}>{\raggedright\arraybackslash}p{4cm}>{\raggedright\arraybackslash}p{6.5cm}}
		\toprule
		\rowcolor{gray!20} \textbf{ID}	& \textbf{Interessato} & \textbf{Task} 
		\\\midrule
	}{
		\bottomrule
	\end{tabular}
}

\newenvironment{changelog}{
	\noindent
	{\Large \textbf{Registro delle modifiche}}
	\noindent
	\begin{table}[h]
		\rowcolors{2}{cyan!80!black!30!}{cyan!80!black!20!}
		\begin{adjustbox}{width=\textwidth}
			\begin{tabular}{|c|c|p{2.35cm}|c|p{3.3cm}|}
				\hline
				\rowcolor{gray!20}
				\textbf{V.} & \textbf{Data} & \textbf{Membro} & \textbf{Ruolo} & \textbf{Descrizione} \\
				
				\hline
			}{
				\hline
			\end{tabular}
		\end{adjustbox}
	\end{table}
	
	\clearpage
}

% 1   2    3        4             5            6            7          
% Ver Data Relatore RuoloRelatore DataVerifica Verificatore Descrizione
\newcommand{\newlog}[7]{
	#1 & #5 & #6 & Verificatore & Controllo qualità \\
	   & #2 & #3 & #4           & #7 \\\hline
}

% INFORMAZIONI DOCUMENTO %%%%%%%%%%%%%%%%%%%%%%%%%
\title{Analisi dei requisiti}
%\date{EMPTY}
\footname{Analisi dei requisiti}

\begin{document}
	% PRIME PAGINE %%%%%%%%%%%%%%%%%%%%%%%%%
	\makefirstpage
	
	% 1   2    3        4             5            6            7          
% Ver Data Relatore RuoloRelatore DataVerifica Verificatore Descrizione
\begin{changelog}
	\newlog{0.3.0}{15/12/2023}{S. Rovea}{Amministratore}{/12/2023}{}{Correzione ed espansione sezione \ref{sec:management_process}}
	\newlog{0.2.0}{02/12/2023}{S. Destro}{Amministratore}{03/12/2023}{A. Feltrin}{Aggiunta sezione su processi organizzativi, sezione \ref{sec:management_process}}
	\newlog{0.1.0}{22/11/2023}{A. Bustreo}{Amministratore}{22/11/2023}{A. Giurisato}{Aggiunto processo supporto per la documentazione, sezione \ref{sec:documentation}. Aggiunta prefazione, sezione \ref{sec:prefazione}}
\end{changelog}
	\clearpage
	
	\tableofcontents
	\clearpage
	
	% CONTENUTO QUI %%%%%%%%%%%%%%%%%%%%%%%%%

	\section{Introduzione}
    	\subsection{Scopo del documento}
    		L'obiettivo di questo documento è quello di raccogliere i casi d’uso ed i requisiti relativi al capitolato proposto da Zextras: JMAP: il nuovo protocollo per la posta elettronica.
  
    
    	\subsection{Scopo del progetto}
    		Il progetto progetto ha come obiettivo lo sviluppo di un servizio di posta elettronica che utilizza il protocollo JMAP.
    
    	\subsection{Glossario}
    		Un glossario in questo documento si rende necessario poiché alcuni termini potrebbero generare dei dubbi sul loro significato.
    		Tutti i termini facenti parte del Glossario verranno contrassegnati con una G.

    	\subsection{Riferimenti}
    		\subsubsection{Riferimenti normativi}
        		\begin{itemize}
            		\item \href{https://www.math.unipd.it/~tullio/IS-1/2023/Progetto/C8.pdf}{Capitolato C8 - JMAP: il nuovo protocollo per la posta elettronica}
        		\end{itemize}

	\section{Descrizione del prodotto}
    	\subsection{Funzioni del prodotto}


\section{Casi d'uso}
    %%%%%%%%%%%%%%% SCENARIO EMAIL %%%%%%%%%%%%%%%%%%%%%%%

    \subsection{UC 1 - Inviare un'email} \label{sec: 1}
    \begin{itemize}
        \item Attore: MUA;
        \item Descrizione: il MUA deve poter scrivere una email;
        \item Scenario principale:
            \begin{enumerate}
            \item il MUA trasmette il destinatario (\hyperref[sec: UC 1.1]{UC 1.1});
            \item il MUA trasmette il mittente (\hyperref[sec: UC 1.2]{UC 1.2});
            \item il MUA trasmette l'oggetto (\hyperref[sec: UC 1.3]{UC 1.3});
            \end{enumerate}
        \item Estensioni: errore... %(\hyperref[sec: UC 1.4.1]{UC 1.4.1});
        \item Precondizioni: l’account che il MUA gestisce è registrato nel sistema ed è autenticato;
        \item Postcondizioni: l'email è stata inviata con successo al destinatario.
    \end{itemize}

    \subsubsection{UC 1.1 - Trasmettere il destinatario} \label{sec: UC 1.1}
    \begin{itemize}
        \item Attore: MUA;
        \item Descrizione: il MUA deve poter trasmettere il destinatario;
        \item Scenario:
        \begin{enumerate}
        \item il MUA si connette con il server;
        \item il MUA trasmette il destinatario.
        \end{enumerate}
        \item Precondizioni: il MUA sta inviando un'email;
        \item Postcondizioni: il destinatario è stato trasmesso.
    \end{itemize}

    \subsubsection{UC 1.2 - Trasmettere il mittente} \label{sec: UC 1.2}
    \begin{itemize}
        \item Attore: MUA;
        \item Descrizione: il MUA deve poter trasmettere il mittente;
        \item Scenario:
        \begin{enumerate}
        \item il MUA si connette con il server;
        \item il MUA trasmette il mittente.
        \end{enumerate}
        \item Estensioni: errore... %(\hyperref[sec: UC 1.4.1]{UC 1.4.1});
        \item Precondizioni: il MUA sta inviando un'email;
        \item Postcondizioni: il mittente è stato trasmesso.
    \end{itemize}
    
    \subsubsection{UC 1.3 - Trasmettere l'oggetto} \label{sec: UC 1.3}
    \begin{itemize}
        \item Attore: MUA;
        \item Descrizione: il MUA deve poter trasmettere il mittente;
        \item Scenario:
        \begin{enumerate}
        \item il MUA si connette con il server;
        \item il MUA trasmette il mittente.
        \end{enumerate}
        \item Estensioni: errore... %(\hyperref[sec: UC 1.4.1]{UC 1.4.1});
        \item Precondizioni: il MUA sta inviando un'email;
        \item Postcondizioni: l'oggetto è stato trasmesso.
    \end{itemize}


    \subsection{UC 2 - Sincronizzare}  \label{sec: UC 2}
    \begin{itemize}
        \item Attore: MUA;
        \item Descrizione: il MUA deve poter aggiornare il suo stato con le nuove informazioni;
        \item Scenario principale:
            \begin{enumerate}
            \item il MUA trasmette una richiesta di aggiornamento (\hyperref[sec: UC 2.1]{UC 2.1}).
            \end{enumerate}
        \item Estensioni: errore... %(\hyperref[sec: UC 1.4.1]{UC 1.4.1});
        \item Precondizioni: l'account che il MUA gestisce è registrato nel sistema ed è autenticato;
        \item Postcondizioni: il MUA è stato aggiornato con le nuove informazioni.
    \end{itemize}

    \subsubsection{UC 2.1 - Trasmettere richiesta di aggiornamento} \label{sec: UC 2.1}
    \begin{itemize}
        \item Attore: MUA;
        \item Descrizione: il MUA deve poter richiedere un aggiornamento dei propri dati;
        \item Scenario:
        \begin{enumerate}
        \item il MUA si connette con il server;
        \item il MUA richiede nuovi dati.
        \end{enumerate}
        \item Estensioni: errore... %(\hyperref[sec: UC 1.4.1]{UC 1.4.1});
        \item Precondizioni: il MUA si sta aggiornando con i nuovi dati;
        \item Postcondizioni: il MUA ha ricevuto nuove informazioni.
    \end{itemize}

    \subsection{UC 3 - Creare un oggetto} \label{sec: UC 3}
    \begin{itemize}
        \item Attore: MUA;
        \item Descrizione: il MUA deve poter creare un oggetto specifico;
        \item Scenario principale:
            \begin{enumerate}
                \item 
            \end{enumerate}
        \item Generalizzazioni:
            \begin{itemize}
                \item il MUA crea una cartella (\hyperref[sec: UC 3.1]{UC 3.1});
                \item il MUA crea un contatto (\hyperref[sec: UC 3.2]{UC 3.2});
                \item il MUA crea una condivisione (\hyperref[sec: UC 3.3]{UC 3.3});
                \item il MUA crea un evento (\hyperref[sec: UC 3.4]{UC 3.4}).
            \end{itemize}
        \item Estensioni: errore ;
        \item Precondizioni: l’account che il MUA gestisce è registrato nel sistema ed è autenticato;
        \item Postcondizioni: è stato creato l’oggetto desiderato.
    \end{itemize}

    \subsubsection{UC 3.1 - Creare una cartella} \label{sec: UC 3.1}
    \begin{itemize}
        \item Attore: MUA;
        \item Descrizione: il MUA deve poter creare una cartella;
        \item Scenario:
            \begin{enumerate}
            \item il MUA si connette con il server;
            \item il MUA trasmette i dati necessari per creare la cartella; 
            \item % scrivere trasmettere solo nome forse???
            \end{enumerate}
        \item Estensioni: Errore nome cartella
        \item Precondizioni: il MUA sta creando una cartella;
        \item Postcondizioni: è stata creata la cartella.
    \end{itemize}

    \subsubsection{UC 3.2 - Creare un contatto} \label{sec: UC 3.2}
    \begin{itemize}
        \item Attore: MUA;
        \item Descrizione: il MUA deve poter creare un contatto;
        \item Scenario:
        \begin{enumerate}
        \item il MUA si connette con il server;
        \item il MUA trasmette i dati necessari per creare il contatto;
        \end{enumerate}
        \item Estensioni: errore...
        \item Precondizioni: il MUA sta creando un contatto;
        \item Postcondizioni: è stata creata il contatto.
    \end{itemize}

    \subsubsection{UC 3.3 - Creare una condivisione} \label{sec: UC 3.3}
    \begin{itemize}
        \item Attore: MUA;
        \item Descrizione: il MUA deve poter creare una condivisione di una cartella;
        \item Scenario:
        \begin{enumerate}
        \item il MUA si connette con il server;
        \item il MUA trasmette i dati necessari per creare la condivisione della cartella.
        \end{enumerate}
        \item Estensioni: errore...
        \item Precondizioni: l’account che il MUA gestisce è registrato nel sistema ed è autenticato;
        \item Postcondizioni: è stata creata la condivisione di una cartella.
    \end{itemize}

    \subsubsection{UC 3.4 - Creare un evento} \label{sec: UC 3.4}
    \begin{itemize}
        \item Attore: MUA;
        \item Descrizione: il MUA deve poter creare un evento;
        \item Scenario:
        \begin{enumerate}
        \item il MUA si connette con il server;
        \item il MUA trasmette i dati necessari per creare un evento.
        \end{enumerate}
        \item Estensioni: errore...
        \item Precondizioni: il MUA sta creando un evento;
        \item Postcondizioni: è stato creato l'evento.
    \end{itemize}

    \subsection{UC 4 - Modificare un oggetto} \label{sec: UC 4}
    \begin{itemize}
        \item Attore: MUA;
        \item Descrizione: il MUA deve poter modificare un oggetto specifico;
        \item Scenario principale:
            \begin{enumerate}
                \item 
            \end{enumerate}
        \item Generalizzazioni:
            \begin{itemize}
            \item il MUA modifica una cartella (\hyperref[sec: UC 4.1]{UC 4.1});
            \item il MUA modifica un contatto (\hyperref[sec: UC 4.2]{UC 4.2});
            \item il MUA mofidica una condivisione (\hyperref[sec: UC 4.3]{UC 4.3});
            \item il MUA modifica un evento (\hyperref[sec: UC 4.4]{UC 4.4}).
            \end{itemize}
        \item Estensioni: errore...
        \item Precondizioni: l’account che il MUA gestisce è registrato nel sistema ed è autenticato;
        \item Postcondizioni: è stato modificato l’oggetto desiderato.
    \end{itemize}

    \subsubsection{UC 4.1 - Modificare una cartella} \label{sec: UC 4.1}
    \begin{itemize}
        \item Attore: MUA;
        \item Descrizione: il MUA deve poter modificare una cartella;
        \item Scenario:
        \begin{enumerate}
        \item il MUA si connette con il server;
        \item Il MUA trasmette i dati necessari per modificare la cartella.
        \end{enumerate}
        \item Estensioni: errore...
        \item Precondizioni: il MUA sta modificando una cartella;
        \item Postcondizioni: è stata modificata la cartella.
    \end{itemize}

    \subsubsection{UC 4.2 - Modificare un contatto} \label{sec: UC 4.2}
    \begin{itemize}
        \item Attore: MUA;
        \item Descrizione: il MUA deve poter modificare un contatto;
        \item Scenario:
        \begin{enumerate}
        \item il MUA si connette con il server;
        \item Il MUA trasmette i dati necessari per modificare il contatto.
        \end{enumerate}
        \item Estensioni: errore...
        \item Precondizioni: il MUA sta modificando un contatto;
        \item Postcondizioni: è stato modificato il contatto.
    \end{itemize}

    \subsubsection{UC 4.3 - Modificare una condivisione} \label{sec: UC 4.3}
    \begin{itemize}
        \item Attore: MUA;
        \item Descrizione: il MUA deve poter modificare una condivisione di una cartella;
        \item Scenario:
        \begin{enumerate}
        \item il MUA si connette con il server;
        \item Il MUA trasmette i dati necessari per modificare la condivisione di una cartella.
        \end{enumerate}
        \item Estensioni: errore...
        \item Precondizioni: il MUA sta modificando una condivisione di una cartella;
        \item Postcondizioni: è stata modificata la condivisione di una cartella.
    \end{itemize}

    \subsubsection{UC 4.4 - Modificare un evento} \label{sec: UC 4.4}
    \begin{itemize}
        \item Attore: MUA;
        \item Descrizione: il MUA deve poter modificare un evento;
        \item Scenario:
        \begin{enumerate}
        \item il MUA si connette con il server;
        \item Il MUA trasmette i dati necessari per modificare un evento.
        \end{enumerate}
        \item Estensioni: 
        \begin{itemize}
            \item Errore dimensioni evento;
        \end{itemize}
        \item Precondizioni: il MUA sta modificando un evento;
        \item Postcondizioni: è stato modificato l'evento.
    \end{itemize}

    \subsection{UC 5 - Eliminare un oggetto} \label{sec: UC 5}
    \begin{itemize}
        \item Attore: MUA;
        \item Descrizione: il MUA deve poter eliminare un oggetto specifico;
        \item Scenario principale:
        \begin{enumerate}
            \item 
        \end{enumerate}
        \item Generalizzazioni:
            \begin{itemize}
            \item il MUA elimina una cartella (\hyperref[sec: UC 5.1]{UC 5.1});
            \item il MUA elimina un contatto (\hyperref[sec: UC 5.2]{UC 5.2});
            \item il MUA elimina una condivisione (\hyperref[sec: UC 5.3]{UC 5.3});
            \item il MUA elimina un evento (\hyperref[sec: UC 5.4]{UC 5.4}).
            \end{itemize}
        \item Estensioni:
        \item Precondizioni: l’account che il MUA gestisce è registrato nel sistema ed è autenticato;
        \item Postcondizioni: è stato eliminato l’oggetto desiderato.
    \end{itemize}

    \subsubsection{UC 5.1 - Eliminare una cartella} \label{sec: UC 5.1}
    \begin{itemize}
        \item Attore: MUA;
        \item Descrizione: il MUA deve poter eliminare una cartella;
        \item Scenario:
        \begin{enumerate}
        \item il MUA si connette con il server;
        \item Il MUA trasmette i dati necessari per eliminare la cartella.
        \end{enumerate}
        \item Estensioni: 
            \begin{itemize}
                \item Errore cartella contenente email;
                \item Errore cartella contenente cartelle;
            \end{itemize}
        \item Precondizioni: il MUA sta eliminando una cartella;
        \item Postcondizioni: è stata eliminata la cartella.
    \end{itemize}

    \subsubsection{UC 5.2 - Eliminare un contatto} \label{sec: UC 5.2}
    \begin{itemize}
        \item Attore: MUA;
        \item Descrizione: il MUA deve poter eliminare un contatto;
        \item Scenario:
        \begin{enumerate}
        \item il MUA si connette con il server;
        \item Il MUA trasmette i dati necessari per eliminare il contatto.
        \end{enumerate}
        \item Estensioni: errore...
        \item Precondizioni: il MUA sta eliminando un contatto;
        \item Postcondizioni: è stato eliminato il contatto.
    \end{itemize}

    \subsubsection{UC 5.3 - Eliminare una condivisione} \label{sec: UC 5.3}
    \begin{itemize}
        \item Attore: MUA;
        \item Descrizione: il MUA deve poter eliminare una condivisione di una cartella;
        \item Scenario:
        \begin{enumerate}
        \item il MUA si connette con il server;
        \item Il MUA trasmette i dati necessari per eliminare la condivisione di una cartella.
        \end{enumerate}
        \item Estensioni: errore...
        \item Precondizioni: il MUA sta eliminando una condivisione di una cartella;
        \item Postcondizioni: è stata eliminata la condivisione di una cartella.
    \end{itemize}

    \subsubsection{UC 5.4 - Eliminare un evento} \label{sec: UC 5.4}
    \begin{itemize}
        \item Attore: MUA;
        \item Descrizione: il MUA deve poter eliminare un evento;
        \item Scenario:
        \begin{enumerate}
        \item il MUA si connette con il server;
        \item Il MUA trasmette i dati necessari per eliminare un evento.
        \end{enumerate}
        \item Estensioni: errore...
        \item Precondizioni: il MUA sta eliminando un evento;
        \item Postcondizioni: è stato eliminato l'evento.
    \end{itemize}


    %creazione cartella problema nome
    \paragraph{UC Errore nome cartella} \label{sec: UC 11.4.2.1}
    \begin{itemize}
        \item Attore: MUA;
        \item Descrizione: il MUA deve ricevere un messaggio di errore adeguato quando si tenta la creazione di una cartella con lo stesso nome e la stessa cartella padre di un'altra;
        \item Scenario:
        \begin{enumerate}
        \item il MUA riceve il messaggio di errore;
        \end{enumerate}   
        \item Precondizioni: 
        \begin{enumerate}
            \item il MUA ha trasmesso il nome per la creazione della cartella;
            \item esiste già una cartella con lo stesso nome e la stessa cartella padre;
        \end{enumerate}
        \item Postcondizioni: il MUA ha ricevuto un messaggio di errore.
    \end{itemize}

    %modificare una cartella (tho anche per creare c'è)
    \paragraph{UC Errore dimensioni cartella} \label{sec: UC 11.4.2.1}
    \begin{itemize}
        \item Attore: MUA;
        \item Descrizione: il MUA deve ricevere un messaggio di errore adeguato quando la modifica di una o più caratteristiche di una cartella superano i limiti di dimensione stabiliti;
        \item Scenario:
        \begin{enumerate}
        \item il MUA riceve il messaggio di errore;
        \end{enumerate}   
        \item Precondizioni: 
        \begin{enumerate}
            \item il MUA ha trasmesso i dati per la modifica di una cartella;
            \item una o più caratteristiche di una cartella superano i limiti di dimensione stabiliti;
        \end{enumerate}
        \item Postcondizioni: il MUA ha ricevuto un messaggio di errore.
    \end{itemize}

    %modificare una cartella (tho anche per creare c'è)
    \paragraph{UC Errore validità cartella} \label{sec: UC 11.4.2.1}
    \begin{itemize}
        \item Attore: MUA;
        \item Descrizione: il MUA deve ricevere un messaggio di errore adeguato quando la modifica di una o più caratteristiche di una cartella riporta un formato non valido;
        \item Scenario:
        \begin{enumerate}
        \item il MUA riceve il messaggio di errore;
        \end{enumerate}   
        \item Precondizioni: 
        \begin{enumerate}
            \item il MUA ha trasmesso i dati per la modifica di una cartella;
            \item una o più caratteristiche di una cartella  non sono valide;
        \end{enumerate}
        \item Postcondizioni: il MUA ha ricevuto un messaggio di errore.
    \end{itemize}


    %modificare una cartella (tho anche per eliminare c'è)
    %troppo? sarebbe una ricerca per id
    \paragraph{UC Errore esistenza cartella} \label{sec: UC 11.4.2.1}
    \begin{itemize}
        \item Attore: MUA;
        \item Descrizione: il MUA deve ricevere un messaggio di errore adeguato quando la ricerca della cartella che si intende modificare fallisce;
        \item Scenario:
        \begin{enumerate}
        \item il MUA riceve il messaggio di errore;
        \end{enumerate}   
        \item Precondizioni: 
        \begin{enumerate}
            \item il MUA ha trasmesso i dati per la modifica di una cartella;
            \item non viene trovata la cartella;
        \end{enumerate}
        \item Postcondizioni: il MUA ha ricevuto un messaggio di errore.
    \end{itemize}

    %eliminazione cartella con email
    \paragraph{UC Errore cartella contenente email} \label{sec: UC 11.4.2.1}
    \begin{itemize}
        \item Attore: MUA;
        \item Descrizione: il MUA deve ricevere un messaggio di errore adeguato quando l'eliminazione di una cartella non è andata a buon fine a causa di una o più email presenti all'interno di quella cartella;
        \item Scenario:
        \begin{enumerate}
        \item il MUA riceve il messaggio di errore;
        \end{enumerate}   
        \item Precondizioni: 
        \begin{enumerate}
            \item il MUA ha trasmesso i dati per l'emilinazione di una cartella ;
            \item la cartella contiene una o più email;
        \end{enumerate}
        \item Postcondizioni: il MUA ha ricevuto un messaggio di errore.
    \end{itemize}

    %eliminazione cartella con cartelle
    \paragraph{UC Errore cartella contenente cartelle} \label{sec: UC 11.4.2.1}
    \begin{itemize}
        \item Attore: MUA;
        \item Descrizione: il MUA deve ricevere un messaggio di errore adeguato quando l'eliminazione di una cartella non è andata a buon fine a causa di una o più cartelle presenti all'interno di quella cartella;
        \item Scenario:
        \begin{enumerate}
        \item il MUA riceve il messaggio di errore;
        \end{enumerate}   
        \item Precondizioni: 
        \begin{enumerate}
            \item il MUA ha trasmesso i dati per l'emilinazione di una cartella ;
            \item la cartella contiene una o più cartelle;
        \end{enumerate}
        \item Postcondizioni: il MUA ha ricevuto un messaggio di errore.
    \end{itemize}

    %modificare un evento/ creazione
    \paragraph{UC Errore dimensioni evento} \label{sec: UC 11.4.2.1}
    \begin{itemize}
        \item Attore: MUA;
        \item Descrizione: il MUA deve ricevere un messaggio di errore adeguato quando la modifica di una o più caratteristiche di un evento superano i limiti di dimensione stabiliti;
        \item Scenario:
        \begin{enumerate}
        \item il MUA riceve il messaggio di errore;
        \end{enumerate}   
        \item Precondizioni: 
        \begin{enumerate}
            \item il MUA ha trasmesso i dati per la modifica di un evento;
            \item una o più caratteristiche di un evento superano i limiti di dimensione stabiliti;
        \end{enumerate}
        \item Postcondizioni: il MUA ha ricevuto un messaggio di errore.
    \end{itemize}

    %modificare un contatto / creazione
    \paragraph{UC Errore dimensioni contatto} \label{sec: UC 11.4.2.1}
    \begin{itemize}
        \item Attore: MUA;
        \item Descrizione: il MUA deve ricevere un messaggio di errore adeguato quando la modifica di una o più caratteristiche di un contatto superano i limiti di dimensione stabiliti;
        \item Scenario:
        \begin{enumerate}
        \item il MUA riceve il messaggio di errore;
        \end{enumerate}   
        \item Precondizioni: 
        \begin{enumerate}
            \item il MUA ha trasmesso i dati per la modifica di un contatto;
            \item una o più caratteristiche di un contatto superano i limiti di dimensione stabiliti;
        \end{enumerate}
        \item Postcondizioni: il MUA ha ricevuto un messaggio di errore.
    \end{itemize}




            %%%% COSE UTILI %%%%
    \paragraph{UC  errore} \label{sec: UC 11.4.2.1}
    \begin{itemize}
        \item Attore: MUA;
        \item Descrizione: il MUA deve ricevere un messaggio di errore adeguato quando l'eliminazione degli elementi non è andata a buon fine;
        \item Scenario:
        \begin{enumerate}
        \item il MUA riceve il messaggio di errore
        \end{enumerate}   
        \item Precondizioni: il MUA ha trasmesso ...;
        \item Postcondizioni: il MUA ha ricevuto un messaggio d'errore.
    \end{itemize}

    \subsection{UC 18 - Stress test}
\begin{itemize}
    \item Attore: developer Zextras;
    \item Descrizione: l'utente deve poter fare degli stess test per misurare le performance del prodotto;
    \item Scenario principale:
        \begin{enumerate}
        \item l'utente avvia uno o una serie di stress test. 
        \end{enumerate}
    \item Precondizioni: l'utente vuole avviare degli stess test;
    \item Postcondizioni: l'utente visualizza i risultati dei test.
\end{itemize}

   

    


%%%%%%%%%%%%%%% SCENARIO EVENTI %%%%%%%%%%%%%%%%%%%%%%%

\subsection{UC 12 - Creazione evento nel calendario}
\begin{itemize}
    \item Attore: utente;
    \item Descrizione: l'utente deve poter creare nuovi eventi nel calendario;
    \item Scenario principale:
        \begin{enumerate}
        \item l'utente inserisce il nome dell'evento (\hyperref[sec: UC 12.1]{UC 12.1});
        \item l'utente inserisce la data (\hyperref[sec: UC 12.2]{UC 12.2});
        \item l'utente inserisce una fascia oraria (\hyperref[sec: UC 12.3]{UC 12.3});
        \item l'utente conferma la creazione dell'evento (\hyperref[sec: UC 12.4]{UC 12.4}).
        \end{enumerate}
    \item Estensioni: l'utente inserisce dei valori non validi nei campi e viene mostrato un messaggio d'errore (\hyperref[sec: UC 12.4.1]{UC 12.4.1});
    \item Precondizioni: l'utente ha acceduto al calendario e vuole creare un nuovo evento;
    \item Postcondizioni: un nuovo evento è stato creato dall'utente.
\end{itemize}

\subsubsection{UC 12.1 - Inserimento nome evento} \label{sec: UC 12.1}
\begin{itemize}
    \item Attore: utente;
    \item Descrizione: l'utente deve poter inserire il nome dell'evento che vuole \par aggiungere al calendario;
    \item Scenario:
        \begin{enumerate}
        \item l'utente seleziona il campo relativo al nome dell'evento;
        \item l'utente digita il nome dell'evento che vuole creare.
        \end{enumerate}
    
    \item Precondizioni: l'utente vuole creare un nuovo evento;
    \item Postcondizioni: l'utente ha compilato il campo relativo al nome dell'evento.
\end{itemize}


\subsubsection{UC 12.2 - Inserimento data evento} \label{sec: UC 12.2}
\begin{itemize}
    \item Attore: utente;
    \item Descrizione: l'utente deve poter inserire la data dell'evento da creare;
    \item Scenario:
        \begin{enumerate}
        \item l'utente seleziona il campo relativo alla data dell'evento;
        \item l'utente inserisce la data per l'evento che vuole creare.
        \end{enumerate}
    
    \item Precondizioni: l'utente vuole creare un nuovo evento;
    \item Postcondizioni: l'utente ha compilato il campo relativo alla data dell'evento.
\end{itemize}


\subsubsection{UC 12.3 - Inserimento ora evento} \label{sec: UC 12.3}
\begin{itemize}
    \item Attore: utente;
    \item Descrizione: l'utente deve poter inserire l'orario dell'evento;
    \item Scenario:
        \begin{enumerate}
        \item l'utente seleziona il campo relativo all'orario dell'evento;
        \item l'utente inserisce la fascia oraria dell'evento che vuole creare.
        \end{enumerate}
    
    \item Precondizioni: l'utente vuole creare un nuovo evento;
    \item Postcondizioni: l'utente ha compilato il campo relativo all'orario dell'evento.
\end{itemize}

\subsubsection{UC 12.4 - Conferma creazione evento} \label{sec: UC 12.4}
\begin{itemize}
    \item Attore: utente;
    \item Descrizione: l'utente deve poter confermare la creazione dell'evento;
    \item Scenario:
        \begin{enumerate}
        \item l'utente conferma la creazione dell'evento.
        \end{enumerate}
    
    \item Precondizioni: l'utente vuole creare un nuovo evento;
    \item Postcondizioni: l'utente ha confermato la creazione dell'evento.
\end{itemize}

\subsubsection{UC 12.4.1 - Visualizzazione errore creazione evento} \label{sec: UC 12.4.1}
\begin{itemize}
    \item Attore: utente;
    \item Descrizione: l'utente deve ricevere un errore a seguito di dati non validi inseriti durante la procedura per creare un evento;
    \item Scenario:
        \begin{enumerate}
        \item l'utente visualizza il messaggio d'errore.
        \end{enumerate}
    
    \item Precondizioni: l'utente vuole creare un nuovo evento e inserisce dei dati non validi;
    \item Postcondizioni: l'utente ha ricevuto un messaggio d'errore.
\end{itemize}



\subsection{UC 13 - Condivisione di un evento}
\begin{itemize}
    \item Attore: utente;
    \item Descrizione: l'utente deve poter condividere degli eventi con altri utenti;
    \item Scenario principale:
        \begin{enumerate}
        \item l'utente sceglie l'evento da condividere (\hyperref[sec: UC 13.1]{UC 13.1});
        \item l'utente inserisce il nome del destinatario dell'evento (\hyperref[sec: UC 13.2]{UC 13.2});
        \item l'utente conferma la condivisione di un evento (\hyperref[sec: UC 13.3]{UC 13.3}).
        \end{enumerate}
    \item Estensioni: l'utente inserisce dei valori non validi nei campi e viene mostrato un messaggio d'errore (\hyperref[sec: UC 13.3.1]{UC 13.3.1});
    \item Precondizioni: l'utente ha acceduto al calendario e vuole condividere un evento;
    \item Postcondizioni: l'utente ha condiviso un evento.
\end{itemize}

\subsubsection{UC 13.1 - Scelta dell'evento da condividere} \label{sec: UC 13.1}
\begin{itemize}
    \item Attore: utente;
    \item Descrizione: l'utente deve scegliere l'evento del calendario che vuole \par condividere;
    \item Scenario:
        \begin{enumerate}
        \item l'utente seleziona l'evento.
        \end{enumerate}
    
    \item Precondizioni: l'utente vuole condividere un evento;
    \item Postcondizioni: l'utente ha scelto l'evento.
\end{itemize}


\subsubsection{UC 13.2 - Scelta del destinatario} \label{sec: UC 13.2}
\begin{itemize}
    \item Attore: utente;
    \item Descrizione: l'utente inserisce il destinatario con cui desidera condividere l'evento;
    \item Scenario:
        \begin{enumerate}
        \item l'utente digita il nome dell'utente destinatario.
        \end{enumerate}
    
    \item Precondizioni: l'utente vuole condividere un evento;
    \item Postcondizioni: l'utente ha scelto il destinatario.
\end{itemize}


\subsubsection{UC 13.3 - Conferma condivisione evento} \label{sec: UC 13.3}
\begin{itemize}
    \item Attore: utente;
    \item Descrizione: l'utente conferma la condivisione di un evento;
    \item Scenario:
        \begin{enumerate}
        \item l'utente deve confermare la condivisione dell'evento scelto con il destinatario.
        \end{enumerate}
    
    \item Precondizioni: l'utente vuole condividere un evento;
    \item Postcondizioni: l'utente confermato la condivisione dell'evento.
\end{itemize}

\subsubsection{UC 13.3.1 - Visualizzazione errore condivisione evento} \label{sec: UC 13.3.1}
\begin{itemize}
    \item Attore: utente;
    \item Descrizione: l'utente deve ricevere un errore a seguito di dati non validi inseriti durante la procedura per condividere un evento;
    \item Scenario:
        \begin{enumerate}
        \item l'utente visualizza il messaggio d'errore.
        \end{enumerate}
    
    \item Precondizioni: l'utente vuole condividere un evento e inserisce dei dati non validi;
    \item Postcondizioni: l'utente ha ricevuto un messaggio d'errore.
\end{itemize}


\subsection{UC 14 - Eliminazione evento nel calendario}
\begin{itemize}
    \item Attore: utente;
    \item Descrizione: l'utente deve poter eliminare eventi nel calendario;
    \item Scenario principale:
        \begin{enumerate}
        \item l'utente sceglie l'evento da eliminare (\hyperref[sec: UC 14.1]{UC 14.1});
        \item l'utente conferma l'eliminazione (\hyperref[sec: UC 14.2]{UC 14.2});
        \item l'utente riceve il messaggio di eliminazione di un evento (\hyperref[sec: UC 14.3]{UC 14.3}).
        \end{enumerate}
    \item Precondizioni: l'utente ha acceduto al calendario e vuole eliminare un evento;
    \item Postcondizioni: l'evento scelto è stato eliminato.
\end{itemize}

\subsubsection{UC 14.1 - Scelta dell'evento da eliminare} \label{sec: UC 14.1}
\begin{itemize}
    \item Attore: utente;
    \item Descrizione: l'utente deve poter scegliere l'evento che vuole eliminare dal calendario;
    \item Scenario:
        \begin{enumerate}
        \item l'utente seleziona l'evento.
        \end{enumerate}
    
    \item Precondizioni: l'utente vuole eliminare un evento;
    \item Postcondizioni: l'utente ha scelto l'evento da eliminare.
\end{itemize}


\subsubsection{UC 14.2 - Conferma eliminazione evento} \label{sec: UC 14.2}
\begin{itemize}
    \item Attore: utente;
    \item Descrizione: l'utente deve confermare l'eliminazione dell'evento scelto;
    \item Scenario:
        \begin{enumerate}
        \item l'utente conferma l'eliminazione.
        \end{enumerate}
    
    \item Precondizioni: l'utente vuole eliminare un evento;
    \item Postcondizioni: l'utente ha confermato l'eliminazione dell'evento.
\end{itemize}


\subsubsection{UC 14.3 - Ricezione messaggio eliminazione} \label{sec: UC 14.3}
\begin{itemize}
    \item Attore: utente;
    \item Descrizione: l'utente deve ricevere un messaggio che garantisce il buon esito dell'operazione;
    \item Scenario:
        \begin{enumerate}
        \item l'utente riceve il messaggio informativo.
        \end{enumerate}
    
    \item Precondizioni: l'utente vuole eliminare un evento;
    \item Postcondizioni: l'utente ha ricevuto il messaggio di corretta eliminazione.
\end{itemize}



%%%%%%%%%%%%%%% SCENARIO CONTATTI %%%%%%%%%%%%%%%%%%%%%%%

\subsection{UC 15 - Creazione nuovo contatto}
\begin{itemize}
    \item Attore: utente;
    \item Descrizione: l'utente deve poter creare nuovi contatti;
    \item Scenario principale:
        \begin{enumerate}
        \item l'utente inserisce il nome del contatto (\hyperref[sec: UC 15.1]{UC 15.1});
        \item l'utente inserisce l'email del nuovo contatto (\hyperref[sec: UC 15.2]{UC 15.2});
        \item l'utente conferma la creazione del contatto (\hyperref[sec: UC 15.3]{UC 15.3}).
        \end{enumerate}
    \item Estensioni: l'utente inserisce dei valori non validi nei campi e viene mostrato un messaggio d'errore (\hyperref[sec: UC 15.3.1]{UC 15.3.1});
    \item Precondizioni: l'utente ha acceduto alla rubrica e vuole creare un nuovo contatto;
    \item Postcondizioni: un nuovo contatto è stato creato dall'utente.
\end{itemize}


\subsubsection{UC 15.1 - Inserimento nome nuovo contatto} \label{sec: UC 15.1}
\begin{itemize}
    \item Attore: utente;
    \item Descrizione: l'utente deve poter inserire il nome del nuovo contatto che vuole aggiungere alla rubrica;
    \item Scenario:
        \begin{enumerate}
        \item l'utente seleziona il campo relativo al nome del contatto;
        \item l'utente digita il nome.
        \end{enumerate}
    
    \item Precondizioni: l'utente vuole creare un nuovo contatto;
    \item Postcondizioni: l'utente ha compilato il campo relativo al nome del contatto.
\end{itemize}


\subsubsection{UC 15.2 - Inserimento email contatto} \label{sec: UC 15.2}
\begin{itemize}
    \item Attore: utente;
    \item Descrizione: l'utente deve poter inserire l'email del nuovo contatto da creare;
    \item Scenario:
        \begin{enumerate}
        \item l'utente seleziona il campo relativo dell'email;
        \item l'utente inserisce l'email del contatto.
        \end{enumerate}
    
    \item Precondizioni: l'utente vuole creare un nuovo contatto;
    \item Postcondizioni: l'utente ha compilato il campo relativo con l'email del contatto.
\end{itemize}


\subsubsection{UC 15.3 - Conferma creazione contatto} \label{sec: UC 15.3}
\begin{itemize}
    \item Attore: utente;
    \item Descrizione: l'utente deve poter confermare la creazione del nuovo contatto;
    \item Scenario:
        \begin{enumerate}
        \item l'utente conferma la creazione.
        \end{enumerate}
    
    \item Precondizioni: l'utente vuole creare un nuovo contatto;
    \item Postcondizioni: l'utente ha confermato la creazione di un nuovo contatto.
\end{itemize}


\subsubsection{UC 15.3.1 - Visualizzazione errore creazione contatto} \label{sec: UC 15.3.1}
\begin{itemize}
    \item Attore: utente;
    \item Descrizione: l'utente deve ricevere un errore a seguito di dati non validi inseriti durante la procedura per creare un nuovo contatto;
    \item Scenario:
        \begin{enumerate}
        \item l'utente visualizza il messaggio d'errore.
        \end{enumerate}
    
    \item Precondizioni: l'utente sta svolgendo la creazione di un contatto e inserisce dei dati non validi;
    \item Postcondizioni: l'utente ha ricevuto un messaggio d'errore.
\end{itemize}


\subsection{UC16 - Condivisione di un contatto}
\begin{itemize}
    \item Attore: utente;
    \item Descrizione: l'utente deve poter condividere un contatto con altri utenti;
    \item Scenario principale:
        \begin{enumerate}
        \item l'utente sceglie il contatto da condividere (\hyperref[sec: UC 16.1]{UC 16.1});
        \item l'utente inserisce il nome del destinatario del utente (\hyperref[sec: UC 16.2]{UC 16.2});
        \item l'utente conferma la condivisione del contatto (\hyperref[sec: UC 16.3]{UC 16.3}).
        \end{enumerate}
    \item Estensioni: l'utente inserisce dei valori non validi nei campi e viene mostrato un messaggio d'errore (\hyperref[sec: UC 16.3.1]{UC 16.3.1});
    \item Precondizioni: l'utente ha acceduto alla rubrica e vuole condividere un contatto;
    \item Postcondizioni: l'utente ha condiviso un contatto.
\end{itemize}

\subsubsection{UC 16.1 - Scelta del contatto da condividere} \label{sec: UC 16.1}
\begin{itemize}
    \item Attore: utente;
    \item Descrizione: l'utente deve scegliere il contatto della rubrica che vuole \par condividere;
    \item Scenario:
        \begin{enumerate}
        \item l'utente seleziona il contatto.
        \end{enumerate}
    
    \item Precondizioni: l'utente vuole condividere un contatto;
    \item Postcondizioni: l'utente ha scelto il contatto.
\end{itemize}


\subsubsection{UC 16.2 - Scelta del destinatario} \label{sec: UC 16.2}
\begin{itemize}
    \item Attore: utente;
    \item Descrizione: l'utente inserisce il destinatario con cui desidera condividere il contatto;
    \item Scenario:
        \begin{enumerate}
        \item l'utente digita il nome dell'utente destinatario.
        \end{enumerate}
    
    \item Precondizioni: l'utente vuole condividere un contatto;
    \item Postcondizioni: l'utente ha scelto il destinatario.
\end{itemize}


\subsubsection{UC 16.3 - Conferma condivisione contatto} \label{sec: UC 16.3}
\begin{itemize}
    \item Attore: utente;
    \item Descrizione: l'utente conferma la condivisione del contatto;
    \item Scenario:
        \begin{enumerate}
        \item l'utente deve confermare la condivisione del contatto scelto con l'utente scelto.
        \end{enumerate}
    
    \item Precondizioni: l'utente vuole condividere un contatto;
    \item Postcondizioni: l'utente ha confermato la condivisione del contatto.
\end{itemize}

\subsubsection{UC 16.3.1 - Visualizzazione errore condivisione contatto} \label{sec: UC 16.3.1}
\begin{itemize}
    \item Attore: utente;
    \item Descrizione: l'utente deve ricevere un errore a seguito di dati non validi inseriti durante la procedura per condividere un contatto;
    \item Scenario:
        \begin{enumerate}
        \item l'utente visualizza il messaggio d'errore.
        \end{enumerate}
    
    \item Precondizioni: l'utente vuole condividere un contatto e inserisce dei dati non validi;
    \item Postcondizioni: l'utente ha ricevuto un messaggio d'errore.
\end{itemize}



\subsection{UC 17 - Eliminazione contatto}
\begin{itemize}
    \item Attore: utente;
    \item Descrizione: l'utente deve poter eliminare un contatto dalla rubrica;
    \item Scenario principale:
        \begin{enumerate}
        \item l'utente sceglie il contatto da eliminare (\hyperref[sec: UC 17.1]{UC 17.1});
        \item l'utente conferma l'eliminazione (\hyperref[sec: UC 17.2]{UC 17.2});
        \item l'utente riceve il messaggio di eliminazione di un contatto (\hyperref[sec: UC 17.3]{UC 17.3}).
        \end{enumerate}
    \item Precondizioni: l'utente ha acceduto alla rubrica e vuole eliminare un contatto;
    \item Postcondizioni: il contatto scelto è stato eliminato.
\end{itemize}

\subsubsection{UC 17.1 - Scelta del contatto da eliminare} \label{sec: UC 17.1}
\begin{itemize}
    \item Attore: utente;
    \item Descrizione: l'utente deve poter scegliere il contatto che vuole eliminare dalla rubrica;
    \item Scenario:
        \begin{enumerate}
        \item l'utente seleziona il contatto.
        \end{enumerate}
    
    \item Precondizioni: l'utente vuole eliminare un contatto;
    \item Postcondizioni: l'utente ha scelto il contatto da eliminare.
\end{itemize}


\subsubsection{UC 17.2 - Conferma eliminazione contatto} \label{sec: UC 17.2}
\begin{itemize}
    \item Attore: utente;
    \item Descrizione: l'utente deve confermare l'eliminazione del contatto scelto;
    \item Scenario:
        \begin{enumerate}
        \item l'utente conferma l'eliminazione.
        \end{enumerate}
    
    \item Precondizioni: l'utente vuole eliminare un contatto;
    \item Postcondizioni: l'utente ha confermato l'eliminazione del contatto.
\end{itemize}


\subsubsection{UC 17.3 - Ricezione messaggio eliminazione} \label{sec: UC 17.3}
\begin{itemize}
    \item Attore: utente;
    \item Descrizione: l'utente deve ricevere un messaggio che garantisce il buon esito dell'operazione;
    \item Scenario:
        \begin{enumerate}
        \item l'utente riceve il messaggio informativo.
        \end{enumerate}
    
    \item Precondizioni: l'utente vuole eliminare un contatto;
    \item Postcondizioni: l'utente ha ricevuto il messaggio di corretta eliminazione.
\end{itemize}

%%%%%%%%%%%%%%% TESTING %%%%%%%%%%%%%%%%%%%%%%

\subsection{UC 18 - Stress test}
\begin{itemize}
    \item Attore: developer Zextras;
    \item Descrizione: l'utente deve poter fare degli stess test per misurare le performance del prodotto;
    \item Scenario principale:
        \begin{enumerate}
        \item l'utente avvia uno o una serie di stress test. 
        \end{enumerate}
    \item Precondizioni: l'utente vuole avviare degli stess test;
    \item Postcondizioni: l'utente visualizza i risultati dei test.
\end{itemize}

\end{document}